% BHU PYQ -- Transcribed & Corrected
% Paper: CHEMJ32: Organic Chemistry-I
% Exam: B.Sc. (Hons.) Semester III (NEP) Examination, 2025-26
% Ref Code: CECONF/N/2025-26/056
% Category: Major
% Department: Chemistry

\documentclass[11pt,a4paper]{article}
\usepackage[a4paper,top=2.5cm,bottom=2.5cm,left=2.8cm,right=2.8cm,headheight=14pt]{geometry}
\usepackage{amsmath,amssymb}
\usepackage[T1]{fontenc}
\usepackage[utf8]{inputenc}
\usepackage{lmodern,microtype,fancyhdr,enumitem,hyperref,lastpage,parskip,listings}

\hypersetup{
    colorlinks=true,
    urlcolor=black,
    linkcolor=black,
    pdftitle={CHEMJ32: Organic Chemistry-I -- BHU Sem III 2025-26 (NEP)}
}

\pagestyle{fancy}
\fancyhf{}
\renewcommand{\headrulewidth}{0.4pt}
\renewcommand{\footrulewidth}{0.4pt}
\fancyhead[L]{\small \textit{Banaras Hindu University}}
\fancyhead[R]{\small \textit{CHEMJ32: Organic Chemistry-I (2025-26)}}
\fancyfoot[C]{\small Page \thepage\ of \pageref{LastPage}}

\newcommand{\pts}[1]{\hfill \textit{[#1]}}

\begin{document}

\begin{center}
  {\large\bfseries BANARAS HINDU UNIVERSITY}\\[2pt]
  {\normalsize B.Sc. (Hons.) Semester III Examination, 2025-26 (NEP)}\\[6pt]
  \rule{0.8\linewidth}{0.5pt}\\[6pt]
  {\large\bfseries CHEMISTRY}\\[4pt]
  {\large\bfseries Paper Code: CHEMJ32 : Organic Chemistry-I}\\[2pt]
  {\small Ref. No.: CECONF/N/2025-26/056}\\[4pt]
  {\normalsize Time: 2$\frac{1}{2}$ Hours\hfill Maximum Marks: 70}
\end{center}

\rule{\linewidth}{0.4pt}

\smallskip

\noindent \textit{Instructions: Answer five questions including Question 1 which is compulsory. The figures in the right-hand margin indicate marks.}

\bigskip

\noindent \textbf{Question 1.}
\begin{enumerate}[label=(\alph*)]
  \item Complete the following reactions with mechanisms: \pts{2 $\times$ 2$\frac{1}{2}$ = 5}
  \begin{enumerate}[label=(\roman*)]
    \item $\text{R-CHO} + \text{Cl-CH}_2\text{-COOEt} \xrightarrow{\text{NaOEt}} ?$
    \item $\text{R-CHO} + \text{Br-CH}_2\text{-COOEt} \xrightarrow[\text{(ii) Hydrolysis}]{\text{(i) Zn, ether}} ?$
  \end{enumerate}
  \item Complete the following reaction with mechanism: \pts{2$\frac{1}{2}$}\\
  $\text{Phenol} \xrightarrow[\text{(ii) H}^+]{\text{(i) CO}_2\text{, NaOH}} ?$
  \item Classify the following molecules (I-IV) as aromatic, anti-aromatic or non-aromatic: \pts{2}
  \begin{enumerate}[label=(\roman*)]
    \item Pyridine / Pyrrole derivative
    \item Boron heterocyclic 5-membered ring
    \item Boron-nitrogen 7-membered ring
    \item Tropone
  \end{enumerate}
  \item Comment on the aromaticity and dipole moment ($\mu$) of the given below compounds I and II: \pts{2$\frac{1}{2}$}
  \begin{enumerate}[label=(\roman*)]
    \item Naphthalene / Azulene derivative
    \item Fulvalene derivative
  \end{enumerate}
  \item Complete the following reaction with mechanism: \pts{2}\\
  $\text{R-COR} + \text{H}_2\text{C}=\text{SMe}_2 \text{ (Sulfur ylide)} \rightarrow ?$
\end{enumerate}

\bigskip

\noindent \textbf{Question 2.}
\begin{enumerate}[label=(\alph*)]
  \item Among the following compounds (I and II) which one have the lowest free energy barrier for cis-trans isomerisation and why? \pts{3}
  \item Draw the Molecular Orbital picture of benzene. \pts{3}
  \item Discuss the mechanism of nitration of benzene with energy profile diagram. Provide evidence in support of its mechanism. \pts{6}
  \item Comment on the nitration of Anisole ($\text{C}_6\text{H}_5\text{-OCH}_3$) using either mixed acid condition ($\text{HNO}_3+\text{H}_2\text{SO}_4$) or $\text{N}_2\text{O}_5\text{ (fuming)}$. \pts{2}
\end{enumerate}

\bigskip

\noindent \textbf{Question 3.}
\begin{enumerate}[label=(\alph*)]
  \item Discuss the Sulfonation of benzene with mechanism. \pts{3}
  \item Discuss the mechanism of Friedel-Crafts Alkylation. Also, point out its limitations. \pts{3+3}
  \item Compare the relative reactivity of benzene, nitrobenzene, toluene and phenol towards electrophilic aromatic substitution reaction. \pts{2}
  \item What would be the orientation (major product formed) when nitration of the following compounds I-VI were performed separately: \pts{3}
  \begin{enumerate}[label=(\roman*)]
    \item Benzoic acid ($\text{Ph-COOH}$)
    \item Toluene ($\text{Ph-CH}_3$)
    \item Nitrobenzene ($\text{Ph-NO}_2$)
    \item Biphenyl
    \item Benzophenone ($\text{Ph-CO-Ph}$)
    \item Phenol ($\text{Ph-OH}$)
  \end{enumerate}
\end{enumerate}

\bigskip

\noindent \textbf{Question 4.}
\begin{enumerate}[label=(\alph*)]
  \item Provide the mechanism of Diazotization of aniline using $\text{HCl}+\text{NaNO}_2$. \pts{4}
  \item How will you prepare m-dibromobenzene starting from benzene? \pts{2}
  \item How will you prepare m-chlorobenzaldehyde starting from benzaldehyde? \pts{2}
  \item How will you prepare 1,3,5-tribromobenzene starting from aniline? \pts{2}
  \item What will be the product when nitrobenzene undergo reduction under neutral condition? \pts{2}
  \item What will be the major products in the following reaction? \pts{2}\\
  $\text{ArN}_2^+\text{Cl}^- + \text{Phenol} \xrightarrow{\text{OH}^-} \text{(A)?} \quad;\quad \text{ArN}_2^+\text{Cl}^- + \text{Aniline} \xrightarrow{\text{H}^+} \text{(B)?}$
\end{enumerate}

\bigskip

\noindent \textbf{Question 5.}
\begin{enumerate}[label=(\alph*)]
  \item Write the synthesis and application of DDT. \pts{3}
  \item What is BHC and how can you synthesize it? \pts{3}
  \item Explain benzyne mechanism with suitable example. Provide evidence in support of its mechanism. \pts{5}
  \item Complete the following reaction with suitable mechanism: \pts{3}\\
  $\text{2-(chloromethyl)pyridine} + \text{PhLi} \rightarrow ?$
\end{enumerate}

\bigskip

\noindent \textbf{Question 6.}
\begin{enumerate}[label=(\alph*)]
  \item Discuss the acid-catalysed Mannich reaction with mechanism. \pts{4}
  \item Discuss the given below Baylis-Hillman reaction with mechanism: \pts{4}\\
  $\text{Ph-CHO} + \text{CH}_2=\text{CH-COOEt} \xrightarrow{\text{DABCO}} ?$
  \item Explain the relative reactivity of acid chloride, acid anhydride, amide and ester towards nucleophilic addition with suitable example. \pts{4}
  \item Complete the following reaction with mechanism: \pts{2}\\
  $\text{Ph-COCH}_2\text{CHO} \xrightarrow{\text{NaOH}} ?$
\end{enumerate}

\bigskip

\noindent \textbf{Question 7.}
\begin{enumerate}[label=(\alph*)]
  \item Among the following compounds (I and II) which one underwent fast solvolysis and why? \pts{2}
  \item Comment on the reactivity of given below compounds (I-III) with acetaldehyde: \pts{2}
  \begin{enumerate}[label=(\roman*)]
    \item $\text{Ph}_3\text{P}^+-\text{CH}_2^-$
    \item $\text{Ph}_3\text{P}^+-\text{CH}^--\text{COR}$
    \item $\text{Ph}_3\text{P}^+-\text{Cyclopentadienyl}^-$
  \end{enumerate}
  \item Discuss the Wittig reaction with mechanism. \pts{3}
  \item Complete the following reaction with suitable mechanism: \pts{2}\\
  $\text{1,5-dicarbonyl compound} \xrightarrow[\text{(ii) H}_2\text{O}]{\text{(i) OH}^-} ?$
  \item Provide the synthetic sequence for the following conversion: \pts{3}\\
  $\text{H}_3\text{C-CHO} \xrightarrow{?} \text{Pentaerythritol } \left[\text{C}(\text{CH}_2\text{OH})_4\right]$
\end{enumerate}

\bigskip

\noindent \textbf{Question 8.} Discuss the mechanism of the following reactions: \pts{4 $\times$ 3$\frac{1}{2}$ = 14}
\begin{enumerate}[label=(\alph*)]
  \item Hell-Volhard-Zelinsky reaction
  \item Perkin reaction
  \item Haloform reaction
  \item Reimer-Tiemann reaction
\end{enumerate}
Complete the reaction: $\text{Ph-COOH} \xrightarrow{\text{SOCl}_2} ? \xrightarrow{\text{CH}_2\text{N}_2} ?$ \pts{2}

\bigskip
\begin{center}
  \textbf{***}
\end{center}

\end{document}
